\par\Needspace{6\baselineskip}\medskip
\SourceLeaf{149}\pdfbookmark[2]{107.}{article-107}
\begin{center}107.\end{center}
Let $\chi$ denote any arc between $-180^0$ and $+180^0$, and $s$ the sine of the arc $\tfrac12\chi$. It is well known that
\[
\tfrac12\chi=s+\tfrac13\cdot\tfrac12s^3+\tfrac15\cdot\frac{1\mathpunct{.}3}{2\mathpunct{.}4}s^5+\tfrac17\cdot\frac{1\mathpunct{.}3\mathpunct{.}5}{2\mathpunct{.}4\mathpunct{.}6}s^7+\text{etc.}
\]
\noindent We also have
\[
\tfrac12\sin\chi=s\BareRoot{(1-ss)}=s-\tfrac12s^3-\frac{1\mathpunct{.}1}{2\mathpunct{.}4}s^5-\frac{1\mathpunct{.}1\mathpunct{.}3}{2\mathpunct{.}4\mathpunct{.}6}s^7-\text{etc.}
\]
\noindent and therefore
\[
\begin{aligned}
\chi-\sin\chi=4\Bigl(&\tfrac13s^3+\tfrac15\cdot\tfrac12s^5\\[.5em]
&+\tfrac17\cdot\frac{1\mathpunct{.}3}{2\mathpunct{.}4}s^7+\tfrac19\cdot\frac{1\mathpunct{.}3\mathpunct{.}5}{2\mathpunct{.}4\mathpunct{.}6}s^9+\text{etc.}\Bigr).
\end{aligned}
\]
\noindent In this series we substitute for $s$ first $\tfrac12\BareRoot{\dfrac{r+r'-\rho}{a}}$ and then $\tfrac12\BareRoot{\dfrac{r+r'+\rho}{a}}$, and multiply the results by $a^{\frac32}$. This gives the following two series, respectively,
\[
\begin{aligned}
&\tfrac16(r+r'-\rho)^{\frac32}+\tfrac1{80}\cdot\frac1a(r+r'-\rho)^{\frac52}\\[.5em]
&\qquad+\tfrac3{1792}\cdot\frac1{aa}(r+r'-\rho)^{\frac72}+\tfrac5{18432}\cdot\frac1{a^3}(r+r'-\rho)^{\frac92}+\text{etc}
\end{aligned}
\]
\[
\begin{aligned}
&\tfrac16(r+r'+\rho)^{\frac32}+\tfrac1{80}\cdot\frac1a(r+r'+\rho)^{\frac52}\\[.5em]
&\qquad+\tfrac3{1792}\cdot\frac1{aa}(r+r'+\rho)^{\frac72}+\tfrac5{18432}\cdot\frac1{a^3}(r+r'+\rho)^{\frac92}+\text{etc.,}
\end{aligned}
\]
\noindent whose sums we denote by $T$ and $U$. Since $2\sin\tfrac12\delta=\pm\BareRoot{\dfrac{r+r'-\rho}{a}}$, with the upper or lower sign according as $2f$ is below or above $180^0$, it is readily seen that $a^{\frac32}(\delta-\sin\delta)=\pm T$, with the sign determined in the same way. Similarly, if we take for $\epsilon$ the smaller value, which is below $180^0$, we have $a^{\frac32}(\epsilon-\sin\epsilon)=U$. Taking the other value, which adds to the first to give $360^0$, clearly gives $a^{\frac32}(\epsilon-\sin\epsilon)=a^{\frac32}\mathpunct{.}360^0-U$. We thus obtain two values for the time $t$
\[
\frac{U\mp T}{k}\qquad\text{and}\qquad\frac{a^{\frac32}\mathpunct{.}360^0}{k}-\frac{U\pm T}{k}\mathpunct{\cdot}
\]
